\section{Ingeniería y validación}

Para validar Compound AI en la ingeniería real, el expositor presentó varios benchmarks: el juego Soang, la búsqueda de rutas en Maze y el diseño de rejillas de 80\u00d780 para óptica a escala nanométrica. En cada categoría se compara un agent (o un latent solver) entrenado con search trace y DPO contra soluciones puras y modelos generativos puros.

\subsection{Puntos de referencia de Soang y Maze}

Soang es un box-pushing game en el que solo se puede empujar hacia adelante: hay que planear bien el orden de las jugadas, porque en cuanto se empuja un box a una esquina ya no se puede regresar. En los experimentos, el Solution-only baseline necesitó cerca de 175 millones de parámetros y muestras de entrenamiento del orden del millón para tener éxito solo de manera ocasional; en cambio, Search Former+DU Forer alcanzó 80\% de exactitud con 15M de parámetros y 100,000 muestras.
El benchmark de Maze usa labyrinth de distintos tamaños: (1) mazes pequeños de sub 20 steps; (2) mapas grandes que se extienden hasta 40+ steps; (3) con obstáculos dinámicos incorporados. El modelo de search trace mostró una generalization más fuerte que el LLM puro, y en los Maze con muchas bifurcaciones, el search trace también ofrece un candidate pruning eficaz.
\begin{knowledgebox}{Resumen de resultados de ingeniería}
En la tarea de Soang, el trace dropout mejoró la robustness; en la tarea de Maze, el consistency check (un judge de autoverificación del agent) redujo de forma significativa la hallucination. Ambos benchmarks muestran que el modelo de trace, con 5 a 10 veces menos datos, logra igualar o incluso superar al solution-only baseline convolucional.
\end{knowledgebox}

\subsection{Benchmark de óptica a escala nanométrica}

En la tarea de diseño óptico, el objetivo es usar un binary pattern de una rejilla de 80\u00d780 para controlar la refracción de cada píxel. El pipeline de entrenamiento requiere: 1) simular la respuesta en frecuencia; 2) calcular el loss del beam splitter y del wavelength multiplexing; 3) tener en cuenta las restricciones del brush durante la fabricación.
El expositor enfatiza: solo incorporando estos indicadores al loss, el latent cost \(C\) puede aprender diseños que sean a la vez válidos y de alto desempeño. Además, el proceso de retroalimentar el solver hacia \(C\) permite que el modelo aprenda de manera automática a partir del failure trace, y así se transfiera entre distintos benchmarks.

\subsection{Resumen de la sección}

Estos benchmarks abarcan tanto el search al estilo del ajedrez como el diseño continuo de óptica, y demuestran que Compound AI puede, en un solo modelo, encadenar múltiples módulos funcionales en distintos dominios y completar el entrenamiento mediante señales de supervisión sencillas.

